\documentclass{article}
\usepackage{amsmath}
\title{Contact-Wrench Audit Mode}
\author{UpstreamDrift Research}
\date{}
\begin{document}
\maketitle
\section{Scope and Conventions}
The existing contact-wrench tracking quadratic program uses joint and root
accelerations and per-foot wrenches in the world frame. Force is in newtons,
moment in newton metres, and the solver scales wrench by body weight $B$.
Its cone slack $s$ is dimensionless, in body-weight-normalized wrench units.
The new opt-in threshold $\epsilon\geq0$ is supplied by the caller; it is a
numerical audit choice, not a physiological or measured force limit.
\section{Policy}
At each control evaluation, strict mode accepts a solution only when support
geometry exists, the solver reports success, all solution arrays have expected
shapes and finite values, slack is nonnegative within solver tolerance, and
$\max_i s_i\leq\epsilon$. Missing support, solver failure, invalid solution,
or excessive slack raises an error rather than returning baseline torque.
The default mode retains baseline fallback; counters expose missing support
and unqualified solves. The pre-existing constraints, model, and controller
physics are unchanged.
\section{Evidence and Limits}
Tests exercise missing support, solver failure, excessive and malformed slack,
a valid bounded result, and default fallback with a real model and a stubbed
solver. They do not establish solver accuracy, applied actuator force,
measured contact, continuous-time constraint satisfaction, or a qualified
swing. The local environment lacks Drake; solver-dependent tests skip rather
than claim verification. Reproduce with the headless focused contact-wrench
pytest suite; set \texttt{QT\_QPA\_PLATFORM=offscreen} and
\texttt{MPLBACKEND=Agg}. This is standalone research documentation.
\end{document}
